% Resume: Arumalla Mohit Krishna
% Compile with pdfLaTeX.

\documentclass[letterpaper,10pt]{article}

\usepackage[margin=0.5in]{geometry}
\usepackage[T1]{fontenc}
\usepackage{lmodern}
\usepackage{microtype}
\usepackage{titlesec}
\usepackage{enumitem}
\usepackage[hidelinks]{hyperref}
\usepackage{tabularx}
\input{glyphtounicode}

\pdfgentounicode=1
\pagestyle{empty}
\urlstyle{same}
\raggedright
\setlength{\parindent}{0pt}
\setlength{\tabcolsep}{0pt}

\titleformat{\section}
  {\large\scshape}{}{0em}{}[\titlerule]

\titlespacing*{\section}{0pt}{6pt}{3pt}

\newcommand{\entry}[2]{%
  \noindent
  \begin{tabularx}{\linewidth}
    {@{}>{\raggedright\arraybackslash}X@{\hspace{6pt}}r@{}}
    \textbf{#1} & #2
  \end{tabularx}\par
}

\newcommand{\bullets}{
  \begin{itemize}[
    leftmargin=1.2em,
    topsep=2pt,
    itemsep=1pt,
    parsep=0pt,
    partopsep=0pt
  ]
}

\newcommand{\finishbullets}{
  \end{itemize}
}

\begin{document}

%----------HEADING----------
\begin{center}
    {\LARGE\bfseries Arumalla Mohit Krishna}\\[3pt]
    +1 (631) 352-8643 $|$
    \href{mailto:mohitkrishna.arumalla@stonybrook.edu}
    {mohitkrishna.arumalla@stonybrook.edu} $|$
    \href{https://linkedin.com/in/mohit-krishna-arumalla}
    {LinkedIn}
\end{center}
\vspace{-8pt}

%-----------EDUCATION-----------
\section{Education}

\entry
  {\textbf{Stony Brook University}}
  {\textbf{Aug 2025 -- Dec 2026}}

M.S. in Computer Science \hfill New York, USA\\
\textit{
Coursework: Natural Language Processing,
Introduction to Computer Vision, Vision Language Modeling
}
\vspace{3pt}

\entry
  {\textbf{Birla Institute of Technology and Science, Pilani}}
  {\textbf{Aug 2019 -- 2023}}

B.E. in Electronics \& Communication; GPA: 8.02/10 \hfill India
%-----------PROFESSIONAL SUMMARY-----------
\section{Professional Summary}

Software engineer with \textbf{2+ years of experience in applied machine learning and data pipeline development}, pursuing an M.S. in Computer Science. Skilled in \textbf{designing, training, and deploying AI/ML models}, building scalable solutions, and integrating data pipelines to meet business and operational goals. Experienced with \textbf{Python, PyTorch, TensorFlow, model evaluation, model retraining}, and cloud infrastructure. Strong analytical, problem-solving, and communication skills to support national security initiatives.

%-----------TECHNICAL SKILLS-----------
\section{Technical Skills}

\textbf{Languages:}
Python, Go, C++, Java, SQL, R\\

\textbf{ML \& Retrieval:}
PyTorch, TensorFlow/Keras, scikit-learn, NumPy, pandas,
Transformers, mBART, BGE, Vision Transformers,
Self-Supervised Learning, LoRA, Visual Prompt Tuning,
Text Embeddings, Vector Search, Model Training and Evaluation, Statistical Analysis\\

\textbf{Data \& Systems:}
gRPC, Kafka, Apache Flink, Flink SQL, Debezium/CDC,
PostgreSQL, Redis, Elasticsearch, Airflow, BigQuery, Data Pipeline Integration\\

\textbf{Infrastructure:}
NVIDIA Triton, AWS, Docker, Kubernetes, Linux, Git,
CI/CD, Istio/Envoy, New Relic, Model Deployment and Monitoring

%-----------EXPERIENCE-----------
\section{Experience}

\entry
  {\textbf{Software Engineer II $|$ Infra.Market, Bangalore}}
  {\textbf{Apr 2024 -- Aug 2025}}

\bullets

  \item Built \textbf{semantic product search with Amazon Titan text embeddings and Elasticsearch vector search}, enabling natural-language query understanding, model serving, and scalable inference pipelines.

  \item Engineered a \textbf{CDC pipeline leveraging Flink, Kafka, Debezium, PostgreSQL, RocksDB, and Elasticsearch} with exactly-once processing semantics to support continuous data flow for ML models with sub-second latency.

  \item Deployed Istio/Envoy across microservices and integrated New Relic distributed tracing to monitor \textbf{model serving infrastructure and maintain system reliability}.

\finishbullets

\entry
  {\textbf{Software Engineer I $|$ Infra.Market, Chennai}}
  {\textbf{Jul 2023 -- Mar 2024}}

\bullets

  \item Owned Search and Location microservices, developing \textbf{gRPC services supporting high-throughput requests} that fed critical data to downstream ML-based ranking and retrieval systems.

  \item Fixed PostgreSQL connection leaks using Gorm ORM, improving resource utilization and service uptime; validated fixes through concurrent load testing at 10K requests/second.

  \item Developed Airflow workflows for BigQuery-to-Elasticsearch data migration with query optimization and automated monitoring supporting \textbf{data governance and pipeline health for ML model inputs}.

\finishbullets

%-----------PROJECTS-----------
\section{Projects}

\entry
  {GPU-Optimized Personalized Product Search}
  {Aug 2026 -- Oct 2026}

\textit{Python, PyTorch, BGE, NVIDIA Triton, Go, gRPC}

\bullets

  \item Developed two-tower retrieval with
  \texttt{BAAI/bge-small-en-v1.5} query/product embeddings
  (384 dimensions) and a trainable user tower combining query
  embeddings with interaction history.

  \item Built an inference pipeline with precomputed product embeddings,
  approximate nearest-neighbor retrieval, and a Go/gRPC gateway;
  integrating Triton for \textbf{GPU-accelerated model serving}.

  \item Designed batching and concurrent-instance benchmarks for
  steady/bursty traffic, measuring p95/p99 latency, throughput,
  GPU utilization, and cost per successful search; evaluated
  Recall@K/NDCG@K while testing overload recovery.

\finishbullets

\entry
  {Attention-Based Hallucination Detection in NMT}
  {Apr 2026 -- May 2026}

\textit{Python, PyTorch, mBART, Transformers, scikit-learn}

\bullets

  \item Built an English-German hallucination-detection pipeline using
  teacher-forced mBART attention features from WMT22/WMT23 quality
  estimation data.

  \item Extracted cross-/self-attention features across 12 decoder
  layers and 16 heads, including entropy, max attention mass,
  variance, top-k mass, and self-to-cross grounding ratios.

  \item Trained token-, sentence-, and head-level classifiers,
  achieving 0.712 F1 and 0.695 ROC-AUC on 50K WMT examples;
  proposed hallucination-aware reranking using localization,
  risk scores, and attention signals.

\finishbullets

\entry
  {Prompt/LoRA for Self-Supervised Learning}
  {Ongoing}

\textit{Python, PyTorch, Vision Transformers}

\bullets

  \item Investigating domain-specific visual prompting and LoRA
  in a self-supervised ViT framework to learn transferable
  representations across image domains.

  \item Studying backbone performance without prompts/LoRA at
  inference; designing comparisons with retained or combined modules
  to evaluate classification performance and parameter efficiency
  under domain shift.

\finishbullets

\end{document}